\documentclass[10pt,a4paper]{amsart}
\usepackage{amsmath,amssymb,amsthm,graphicx}
\usepackage{verbatim}
\newenvironment{reference}{\comment}{\endcomment}
%\newenvironment{reference}{}{}
\newenvironment{slogan}{\comment}{\endcomment}
\newenvironment{history}{\comment}{\endcomment}
% For commutative diagrams we use Xy-pic
\usepackage[all]{xy}

% We use 2cell for 2-commutative diagrams.
\xyoption{2cell}
% Theorem environments.
%
\theoremstyle{plain}
\newtheorem{theorem}[subsection]{Theorem}
\newtheorem{proposition}[subsection]{Proposition}
\newtheorem{lemma}[subsection]{Lemma}

\theoremstyle{definition}
\newtheorem{definition}[subsection]{Definition}
\newtheorem{example}[subsection]{Example}
\newtheorem{exercise}[subsection]{Exercise}
\newtheorem{situation}[subsection]{Situation}

\theoremstyle{remark}
\newtheorem{remark}[subsection]{Remark}
\newtheorem{remarks}[subsection]{Remarks}

\numberwithin{equation}{subsection}

% Macros
%
\def\lim{\mathop{\mathrm{lim}}\nolimits}
\def\colim{\mathop{\mathrm{colim}}\nolimits}
\def\Spec{\mathop{\mathrm{Spec}}}
\def\Hom{\mathop{\mathrm{Hom}}\nolimits}
\def\Ext{\mathop{\mathrm{Ext}}\nolimits}
\def\SheafHom{\mathop{\mathcal{H}\!\mathit{om}}\nolimits}
\def\SheafExt{\mathop{\mathcal{E}\!\mathit{xt}}\nolimits}
\def\Sch{\mathit{Sch}}
\def\Mor{\mathop{\mathrm{Mor}}\nolimits}
\def\Ob{\mathop{\mathrm{Ob}}\nolimits}
\def\Sh{\mathop{\mathit{Sh}}\nolimits}
\def\NL{\mathop{N\!L}\nolimits}
\def\CH{\mathop{\mathrm{CH}}\nolimits}
\def\proetale{{pro\text{-}\acute{e}tale}}
\def\etale{{\acute{e}tale}}
\def\QCoh{\mathit{QCoh}}
\def\Ker{\mathop{\mathrm{Ker}}}
\def\Im{\mathop{\mathrm{Im}}}
\def\Coker{\mathop{\mathrm{Coker}}}
\def\Coim{\mathop{\mathrm{Coim}}}

% Boxtimes
%
\DeclareMathSymbol{\boxtimes}{\mathbin}{AMSa}{"02}

%
% Macros for moduli stacks/spaces
%
\def\QCohstack{\mathcal{QC}\!\mathit{oh}}
\def\Cohstack{\mathcal{C}\!\mathit{oh}}
\def\Spacesstack{\mathcal{S}\!\mathit{paces}}
\def\Quotfunctor{\mathrm{Quot}}
\def\Hilbfunctor{\mathrm{Hilb}}
\def\Curvesstack{\mathcal{C}\!\mathit{urves}}
\def\Polarizedstack{\mathcal{P}\!\mathit{olarized}}
\def\Complexesstack{\mathcal{C}\!\mathit{omplexes}}
% \Pic is the operator that assigns to X its picard group, usage \Pic(X)
% \Picardstack_{X/B} denotes the Picard stack of X over B
% \Picardfunctor_{X/B} denotes the Picard functor of X over B
\def\Pic{\mathop{\mathrm{Pic}}\nolimits}
\def\Picardstack{\mathcal{P}\!\mathit{ic}}
\def\Picardfunctor{\mathrm{Pic}}
\def\Deformationcategory{\mathcal{D}\!\mathit{ef}}
\usepackage{cleveref}
\crefname{originalreference}{original reference}{original references}
\Crefname{originalreference}{Original reference}{Original references}
\makeatletter
\newlabel{algebra-lemma-depth-ext}{{algebra-lemma-depth-ext}{source}{}{sourcepointer-1}{}}
\newlabel{algebra-lemma-depth-ext@cref}{{[originalreference][0][]algebra-lemma-depth-ext}{[1][0][]source}{}{}{}}
\makeatother
\crefname{originallemma}{original lemma}{original lemmas}
\Crefname{originallemma}{Original lemma}{Original lemmas}
\makeatletter
\newlabel{lemma-depth}{{lemma-depth}{source}{}{sourcepointer-2}{}}
\newlabel{lemma-depth@cref}{{[originallemma][0][]lemma-depth}{[1][0][]source}{}{}{}}
\makeatother
\crefname{originallemma}{original lemma}{original lemmas}
\Crefname{originallemma}{Original lemma}{Original lemmas}
\makeatletter
\newlabel{lemma-dualizing}{{lemma-dualizing}{source}{}{sourcepointer-3}{}}
\newlabel{lemma-dualizing@cref}{{[originallemma][0][]lemma-dualizing}{[1][0][]source}{}{}{}}
\makeatother
\crefname{originallemma}{original lemma}{original lemmas}
\Crefname{originallemma}{Original lemma}{Original lemmas}
\makeatletter
\newlabel{lemma-sitting-in-degrees}{{lemma-sitting-in-degrees}{source}{}{sourcepointer-4}{}}
\newlabel{lemma-sitting-in-degrees@cref}{{[originallemma][0][]lemma-sitting-in-degrees}{[1][0][]source}{}{}{}}
\makeatother
\crefname{originallemma}{original lemma}{original lemmas}
\Crefname{originallemma}{Original lemma}{Original lemmas}
\makeatletter
\newlabel{lemma-special-case-local-duality}{{lemma-special-case-local-duality}{source}{}{sourcepointer-5}{}}
\newlabel{lemma-special-case-local-duality@cref}{{[originallemma][0][]lemma-special-case-local-duality}{[1][0][]source}{}{}{}}
\makeatother
\title{The depth as the smallest degree in which a normalized dualizing complex is nonzero}
\author{The Stacks Project Authors}
\date{stacks-project:dualizing.tex:local-snapshot-20261012; GNU FDL 1.2 or later}
\begin{document}
\maketitle
\noindent\textit{Editorial note.} All delivered mathematical text is the author's own, copied from the pinned Stacks Project chapter source sources/stacks-dualizing/pinned-original-source.tex; no statement or proof step was rewritten, reordered or removed.\par\medskip
\noindent\textit{Editorial note.} Preamble provenance: the Stacks Project keeps one single shared preamble for every chapter; this card takes it directly from the staged sources/stacks-dualizing/preamble.tex as three verbatim spans (lines 8-12, 14-18 and 154-217), so no derived file was written and nothing was added or redefined. The wrapper class and the base packages are supplied by the assembler.\par\medskip
\noindent\textit{Editorial note.} Bibliography provenance: neither the counted statement nor its proof contains a citation (this lane rescans the whole span before writing the card), so the card carries no bibliography and no entry was reconstructed.\par\medskip
\noindent\textit{Editorial note.} Selection provenance: this card was opened under the file-based control contract (CONTROL.md), with the mechanical pre-card gate run on its statement and proof interval before the card was written.\par\medskip
\noindent\textit{Allowed standard prerequisites (editorial).}\begin{itemize}
\item The degree confinement for the cohomology of a normalized dualizing complex against a finite module. Invoked at the start of the counted proof (source line 3817). Gives the result directly, and is the reference point for the two alternative arguments.
\item The duality statement identifying Ext of a module into the dualizing complex with Ext of the residue field into the ring. Invoked inside the first alternative (source lines 3827-3829). Rewrites the smallest nonvanishing degree as the smallest degree in which Ext of the residue field into the ring is nonzero.
\item The characterization of the depth as the smallest degree in which Ext of the residue field into the ring is nonzero. Invoked inside the first alternative (source lines 3830-3831). Identifies that smallest degree with the depth.
\item The vanishing criterion for local cohomology in degree equal to the depth. Invoked inside the second alternative (source lines 3836-3838). Gives the second derivation of the statement through local duality.
\end{itemize}\medskip
\ifdefined\hypertarget\hypertarget{sourcepointer-1}{}\fi
\noindent\textit{Source reference.} Original reference algebra-lemma-depth-ext: sources/stacks-dualizing/pinned-original-source.tex, source lines 3831--3831. Navigation pointer only: the counted text refers to the Stacks item algebra-lemma-depth-ext, which belongs to another chapter of the project and is not part of this card. The pointer carries the target's own label as its number because this chapter's staged snapshot does not contain the other chapter's tag table; its content is not reproduced here.\par\medskip
\ifdefined\hypertarget\hypertarget{sourcepointer-2}{}\fi
\noindent\textit{Source reference.} Original lemma lemma-depth: sources/stacks-dualizing/pinned-original-source.tex, source lines 1921--1921. Navigation pointer only: the counted text refers to the Stacks item lemma-depth, which is an existing labeled result of this same chapter and is not part of this card. The pointer carries the target's own label as its number and its definition line as its span; its content is not reproduced here.\par\medskip
\ifdefined\hypertarget\hypertarget{sourcepointer-3}{}\fi
\noindent\textit{Source reference.} Original lemma lemma-dualizing: sources/stacks-dualizing/pinned-original-source.tex, source lines 2818--2818. Navigation pointer only: the counted text refers to the Stacks item lemma-dualizing, which is an existing labeled result of this same chapter and is not part of this card. The pointer carries the target's own label as its number and its definition line as its span; its content is not reproduced here.\par\medskip
\ifdefined\hypertarget\hypertarget{sourcepointer-4}{}\fi
\noindent\textit{Source reference.} Original lemma lemma-sitting-in-degrees: sources/stacks-dualizing/pinned-original-source.tex, source lines 3232--3232. Navigation pointer only: the counted text refers to the Stacks item lemma-sitting-in-degrees, which is an existing labeled result of this same chapter and is not part of this card. The pointer carries the target's own label as its number and its definition line as its span; its content is not reproduced here.\par\medskip
\ifdefined\hypertarget\hypertarget{sourcepointer-5}{}\fi
\noindent\textit{Source reference.} Original lemma lemma-special-case-local-duality: sources/stacks-dualizing/pinned-original-source.tex, source lines 3727--3727. Navigation pointer only: the counted text refers to the Stacks item lemma-special-case-local-duality, which is an existing labeled result of this same chapter and is not part of this card. The pointer carries the target's own label as its number and its definition line as its span; its content is not reproduced here.\par\medskip
\begin{lemma}
\label{lemma-depth-in-terms-dualizing-complex}
Let $(A, \mathfrak m, \kappa)$ be a Noetherian local ring with
normalized dualizing complex $\omega_A^\bullet$.
Then $\text{depth}(A)$ is equal to the smallest integer $\delta \geq 0$
such that $H^{-\delta}(\omega_A^\bullet) \not = 0$.
\end{lemma}

\begin{proof}
This follows immediately from
Lemma \ref{lemma-sitting-in-degrees}.
Here are two other ways to see that it is true.

\medskip\noindent
First alternative. By Nakayama's lemma we see that
$\delta$ is the smallest integer such that
$\Hom_A(H^{-\delta}(\omega_A^\bullet), \kappa) \not = 0$.
In other words, it is the smallest integer such that
$\Ext_A^{-\delta}(\omega_A^\bullet, \kappa)$
is nonzero. Using Lemma \ref{lemma-dualizing} and the fact that
$\omega_A^\bullet$ is normalized this is equal to the
smallest integer such that $\Ext_A^\delta(\kappa, A)$ is
nonzero. This is equal to the depth of $A$ by
Algebra, Lemma \ref{algebra-lemma-depth-ext}.

\medskip\noindent
Second alternative. By the local duality theorem
(in the form of Lemma \ref{lemma-special-case-local-duality})
$\delta$ is the smallest integer such that $H^\delta_\mathfrak m(A)$
is nonzero. This is equal to the depth of $A$ by
Lemma \ref{lemma-depth}.
\end{proof}

\end{document}
